% !TEX program = xelatex
\documentclass[12pt,a4paper]{article}

\usepackage[a4paper,margin=2.5cm]{geometry}
\usepackage{graphicx}
\graphicspath{{figures/}}

\usepackage{tikz}
\usetikzlibrary{positioning,arrows.meta,calc,fit,shapes.geometric}

% ---- Tiếng Việt + Unicode ----
\usepackage{fontspec}
\usepackage{polyglossia}
\setmainlanguage{vietnamese}
\setotherlanguage{english}
\setmainfont{TeX Gyre Termes}
\setsansfont{TeX Gyre Heros}
\setmonofont{TeX Gyre Cursor}
\newcommand{\scenario}[1]{\texttt{#1}}

\usepackage{tabularx}
\usepackage{microtype}
\usepackage{booktabs}
\usepackage{multirow}
\usepackage{array}
\usepackage{float}
\usepackage[font=small,labelfont=bf]{caption}
\usepackage{algorithm}
\usepackage{algpseudocode}
\usepackage{enumitem}
\usepackage{amsmath,amssymb,amsfonts,mathtools}
\usepackage{hyperref}
\hypersetup{
  colorlinks=true,
  linkcolor=black,
  urlcolor=blue,
  citecolor=black
}

\numberwithin{equation}{section}

\setlength{\parskip}{0.35em}
\setlength{\parindent}{1em}
\renewcommand{\arraystretch}{1.15}
\tolerance=2000
\emergencystretch=3em

\begin{document}

% ---- Nhãn tiếng Việt ----
\renewcommand{\contentsname}{MỤC LỤC}
\renewcommand{\listfigurename}{DANH MỤC HÌNH VẼ}
\renewcommand{\listtablename}{DANH MỤC BẢNG BIỂU}
\renewcommand{\refname}{TÀI LIỆU THAM KHẢO}
\renewcommand{\figurename}{Hình}
\renewcommand{\tablename}{Bảng}

% =====================================================================
%  TRANG BÌA
% =====================================================================
\begin{titlepage}
    \centering
    \setlength\parindent{0cm}

    \vspace*{-1cm}
    \begin{tikzpicture}[overlay, remember picture]
        \draw[line width=3pt,color=black,fill=none]
            ($(current page.north west)+(2.5cm,-2cm)$) rectangle ($(current page.south east)-(1.5cm,-2.5cm)$);
        \draw[line width=1pt,color=black,fill=none]
            ($(current page.north west)+(2.65cm,-2.15cm)$) rectangle ($(current page.south east)-(1.65cm,-2.65cm)$);
    \end{tikzpicture}

    \fontsize{14pt}{16pt}\selectfont\textbf{HỌC VIỆN CÔNG NGHỆ BƯU CHÍNH VIỄN THÔNG}\par
    \vspace{0.3cm}
    \underline{\textbf{KHOA SAU ĐẠI HỌC}}\par

    \vspace{1cm}

    \includegraphics[scale=0.25]{logo-PTIT.jpg}\par

    \vspace{1.4cm}
    \fontsize{20pt}{24pt}\selectfont\textbf{\MakeUppercase{Báo cáo môn học}}\par

    \vspace{1cm}
    \fontsize{18pt}{21pt}\selectfont\textbf{Bám đối tượng bằng MeanShift / SIFT}\par

    \vspace{0.4cm}
    \fontsize{13pt}{16pt}\selectfont\textit{(Object tracking using Mean Shift / Scale-Invariant Feature Transform)}\par

    \vspace{1.2cm}
    \begin{center}
        \begin{minipage}{0.88\textwidth}
        \fontsize{12pt}{15pt}\selectfont
        \renewcommand{\arraystretch}{1.3}
        \begin{tabularx}{\textwidth}{@{}>{\bfseries}l X@{}}
        Giảng viên hướng dẫn: & \textbf{PGS.TS. PHẠM VĂN CƯỜNG} \\[0.8em]
        Nhóm thực hiện: & \textbf{3} \\[0.8em]
        Lớp: & \textbf{M25CQHT02-B} \\[0.8em]
        Học viên thực hiện: &
        \textbf{NGUYỄN MẠNH CƯỜNG} \quad B25CHHT082 \\
        & \textbf{TRẦN DUY KHÁNH} \quad B25CHHT100 \\
        & \textbf{NGUYỄN ĐỨC HOÀNG} \quad B25CHHT096 \\
        & \textbf{HÀ PHƯƠNG NGUYÊN} \quad B25CHHT108 \\
        & \textbf{TRẦN ĐÌNH DŨNG} \quad B25CHHT089 \\
        & \textbf{PHẠM MINH HIẾU} \quad B25CHHT095 \\
        \end{tabularx}
        \end{minipage}
    \end{center}

    \vfill
    \fontsize{12pt}{14pt}\selectfont\textbf{HÀ NỘI -- 2026}\par
    \vspace{0.5cm}
\end{titlepage}

% =====================================================================
% =====================================================================
%  PHẦN ĐẦU
% =====================================================================
\pagenumbering{roman}

\tableofcontents
\clearpage

\listoffigures
\vspace{1.5em}
\listoftables
\clearpage

\section*{DANH MỤC TỪ VIẾT TẮT}
\addcontentsline{toc}{section}{DANH MỤC TỪ VIẾT TẮT}

\begin{table}[H]
\centering
\small
\begin{tabularx}{\textwidth}{@{}l l X@{}}
\toprule
\textbf{Viết tắt} & \textbf{Nguyên gốc} & \textbf{Nghĩa tiếng Việt} \\
\midrule
AUC & Area Under the Curve & Diện tích dưới đường cong \\
BWH & Background-Weighted Histogram & Histogram trọng số nền \\
CAMShift & Continuously Adaptive Mean Shift & Mean shift thích nghi liên tục \\
CBWH & Corrected Background-Weighted Histogram & Histogram trọng số nền hiệu chỉnh \\
CLE & Center Location Error & Sai số vị trí tâm \\
DoG & Difference of Gaussians & Sai khác Gauss \\
DSST & Discriminative Scale Space Tracker & Bộ bám không gian tỉ lệ phân biệt \\
EM & Expectation--Maximization & Kỳ vọng và Cực đại hoá \\
FPS & Frames Per Second & Số khung hình mỗi giây \\
IoU & Intersection over Union & Tỉ lệ giao trên hợp \\
KDE & Kernel Density Estimation & Ước lượng mật độ hạt nhân \\
LBP & Local Binary Pattern & Mẫu nhị phân cục bộ \\
ORB & Oriented FAST and Rotated BRIEF & Giữ nguyên tên gốc \\
OTB & Object Tracking Benchmark & Bộ chuẩn bám đối tượng \\
RANSAC & RANdom SAmple Consensus & Đồng thuận mẫu ngẫu nhiên \\
SIFT & Scale-Invariant Feature Transform & Biến đổi đặc trưng bất biến tỉ lệ \\
SOAMST & Scale and Orientation Adaptive Mean Shift Tracking & Mean shift thích nghi tỉ lệ và hướng \\
SOT & Single Object Tracking & Bám đối tượng đơn \\
SURF & Speeded Up Robust Features & Đặc trưng bền vững tăng tốc \\
\bottomrule
\end{tabularx}
\end{table}
\clearpage

% =====================================================================
\section*{LỜI MỞ ĐẦU}
\addcontentsline{toc}{section}{LỜI MỞ ĐẦU}

Bám đối tượng thị giác là bài toán nền tảng của thị giác máy tính, giữ vai trò khối xử lý trung
gian giữa tầng tri giác cấp thấp và tầng suy luận ngữ nghĩa cấp cao trong hầu hết các hệ thống
phân tích video.

Báo cáo tập trung vào ba nhóm có quan hệ trực tiếp: họ thuật toán MeanShift cùng các biến thể,
họ đặc trưng cục bộ với SIFT làm trọng tâm, và các phương án lai kết hợp hai hướng trên. Trong phần
thực nghiệm, KCF được bổ sung như một baseline cổ điển thuộc họ bộ lọc tương quan để cung cấp mốc
tham chiếu ngoài hai cách biểu diễn chính. KCF chỉ được trình bày ở mức cần thiết cho thiết kế thực
nghiệm; các phương pháp học sâu vẫn nằm ngoài phạm vi báo cáo.

Về đối chiếu với yêu cầu của học phần, cấu trúc bảy mục của báo cáo tương ứng như sau: Mục~1 trình
bày \textit{bài toán}; các Mục~2, 3, 4 và 5 trình bày \textit{phương pháp}; Mục~6 tổng hợp
\textit{kết quả thực nghiệm} đã công bố; Mục~7 là phần \textit{thảo luận và kết luận}.

Bên cạnh phần tổng hợp lý thuyết, nhóm thực hiện một thực nghiệm đối chứng có khả năng tái lập trên
cùng máy, cùng dữ liệu và cùng cách khởi tạo. Bảy bộ bám được so sánh gồm MeanShift,
CBWH-MeanShift, CAMShift, SOAMST, SIFT, MeanShift+SIFT và KCF. Mục tiêu của thực nghiệm không phải
chứng minh một thuật toán mới, mà kiểm tra trực tiếp các nhận định về ảnh hưởng của mô hình nền,
khả năng thích nghi hình học bằng CAMShift/SOAMST, độ bền của đặc trưng SIFT, lợi ích của phương án
lai và vị trí của một baseline bộ lọc tương quan.
\clearpage

\pagenumbering{arabic}
\setcounter{page}{1}

% =====================================================================
\section{BÀI TOÁN BÁM ĐỐI TƯỢNG}
% =====================================================================

\subsection{Đầu vào và đầu ra}

Bám đối tượng đơn (Single Object Tracking, viết tắt là SOT) được phát biểu theo thiết lập trực tuyến, khởi tạo từ hộp bao ở khung hình đầu~\cite{smeulders2014}.

\textbf{Đầu vào} gồm hai thành phần: một chuỗi video $I_1, I_2, \dots, I_T$ và trạng thái ban đầu
$\mathbf{s}_1 = (x_1, y_1, w_1, h_1)$ của đối tượng tại khung hình đầu tiên, thường là một hộp bao
do người dùng đánh dấu.

\textbf{Đầu ra} là dãy trạng thái ước lượng $\hat{\mathbf{s}}_2, \hat{\mathbf{s}}_3, \dots,
\hat{\mathbf{s}}_T$ của đối tượng ở các khung hình còn lại, thu được từ hàm ước lượng
\begin{equation}
\hat{\mathbf{s}}_t = \mathcal{F}\big(I_t;\; I_{1:t-1},\, \mathbf{s}_1,\, \hat{\mathbf{s}}_{1:t-1}\big),
\qquad t = 2, 3, \dots, T .
\end{equation}

Hai đặc điểm sau quy định bản chất của bài toán. Thứ nhất, bài toán \textit{không phụ thuộc mô
hình lớp}: đối tượng chỉ được xác định bởi chú thích ở khung hình đầu, hệ thống không có sẵn mô
hình lớp huấn luyện trước. Đây là điểm phân biệt căn bản giữa bám đối tượng và phát hiện đối
tượng. Thứ hai, bài toán mang tính \textit{trực tuyến và nhân quả}: tại thời điểm $t$ chỉ được sử
dụng thông tin từ các khung hình quá khứ.

Bảng~\ref{tab:detection-tracking} làm rõ ranh giới giữa phát hiện và bám đối tượng. Phát hiện trả
lời câu hỏi ``trong ảnh có những đối tượng nào?'', còn bám đối tượng duy trì trạng thái của
\textit{một mục tiêu đã biết} qua thời gian.

\begin{table}[H]
\centering
\caption{Phân biệt phát hiện đối tượng và bám đối tượng đơn}
\label{tab:detection-tracking}
\small
\begin{tabular}{@{}p{3.0cm}p{5.1cm}p{5.4cm}@{}}
\toprule
\textbf{Tiêu chí} & \textbf{Phát hiện đối tượng} & \textbf{Bám đối tượng đơn} \\
\midrule
Đầu vào & Một ảnh hoặc từng khung hình được xử lý độc lập & Chuỗi video và hộp bao mục tiêu ở khung đầu \\
Đầu ra & Hộp bao và nhãn của các đối tượng trong mỗi ảnh & Dãy trạng thái của cùng một mục tiêu qua các khung hình \\
Thông tin thời gian & Không bắt buộc khai thác lịch sử & Dùng kết quả quá khứ để cập nhật trạng thái hiện tại \\
Mô hình mục tiêu & Thường dựa trên mô hình lớp đã huấn luyện & Mục tiêu được xác định trực tiếp bởi chú thích ban đầu \\
Mục tiêu chính & Nhận biết và định vị các thể hiện trong ảnh & Duy trì danh tính, vị trí và kích thước của mục tiêu \\
\bottomrule
\end{tabular}
\end{table}

Từ hai đặc điểm trên nảy sinh khó khăn trung tâm của bài toán: ngoại hình đối tượng biến đổi liên
tục theo thời gian, trong khi mô hình biểu diễn lại buộc phải tự cập nhật từ chính kết quả ước
lượng của mình. Vòng phản hồi khép kín này dẫn tới hiện tượng \textit{trôi} (drift) và tích luỹ
sai số.

\subsection{Các thách thức chính}

Bảng~\ref{tab:challenges} chọn các thuộc tính thách thức có liên quan trực tiếp tới MeanShift và SIFT từ hệ phân loại của Object Tracking Benchmark (OTB)~\cite{wu2015}.

\begin{table}[H]
\centering
\caption{Các thách thức chính của bài toán bám đối tượng}
\label{tab:challenges}
\small
\begin{tabular}{@{}p{4.0cm}p{5.3cm}p{5.2cm}@{}}
\toprule
\textbf{Thách thức} & \textbf{Biểu hiện} & \textbf{Ảnh hưởng tới MeanShift/SIFT} \\
\midrule
Thay đổi ánh sáng & Độ sáng vùng đối tượng biến thiên & Histogram màu bị lệch, MeanShift dễ trôi \\
Thay đổi tỉ lệ & Đối tượng phóng to hoặc thu nhỏ & MeanShift bản nguyên với cửa sổ cố định thất bại \\
Phép quay & Đối tượng xoay trong hoặc ngoài mặt phẳng ảnh & SIFT xử lý tốt, MeanShift cần biến thể SOAMST \\
Che khuất & Che một phần hoặc toàn phần & MeanShift chịu được che một phần, mất mục tiêu khi che toàn phần \\
Nền giống đối tượng & Nền chứa thành phần màu trùng đối tượng & Điểm yếu căn bản nhất của MeanShift \\
Chuyển động nhanh & Dịch chuyển lớn giữa hai khung hình liên tiếp & Vùng tìm kiếm cục bộ không phủ được \\
Đối tượng ít texture & Bề mặt trơn, không có chi tiết & Không đủ số điểm đặc trưng cho SIFT \\
\bottomrule
\end{tabular}
\end{table}

Có thể nhận thấy các thách thức trong Bảng~\ref{tab:challenges} phân bố khá đối xứng giữa hai
phương pháp: những tình huống gây khó cho MeanShift lại thường là thế mạnh của SIFT và ngược lại.
Quan sát này là cơ sở cho các phương án lai được trình bày ở Mục~5.

\subsection{Phân loại các hướng tiếp cận}

Trong phạm vi báo cáo, các phương pháp bám được phân thành ba nhóm theo cách biểu diễn đối tượng.

\textbf{Bám dựa trên histogram} biểu diễn đối tượng bằng phân bố màu toàn cục của vùng ảnh và định
vị bằng cách tìm vùng có phân bố tương đồng nhất. MeanShift và các biến thể thuộc nhóm này.

\textbf{Bám dựa trên đặc trưng cục bộ} biểu diễn đối tượng bằng một tập điểm đặc trưng kèm vector
mô tả, sau đó bám bằng cách đối sánh điểm giữa các khung hình và ước lượng phép biến đổi hình học.
SIFT, SURF và ORB thuộc nhóm này.

\textbf{Bám lai đa đặc trưng} kết hợp hai cách biểu diễn trên nhằm tận dụng ưu điểm bổ trợ của
chúng.

Ngoài ba nhóm trên còn có hướng bám bằng bộ lọc tương quan và hướng học sâu, nhưng hai hướng này
nằm ngoài phạm vi báo cáo.

% =====================================================================
\section{HỌ THUẬT TOÁN MEANSHIFT}
% =====================================================================

\subsection{MeanShift nguyên bản}

Phần này trình bày mô hình bám dựa trên hạt nhân của Comaniciu, Ramesh và Meer~\cite{comaniciu2003}; trích dẫn được đặt tại đây để bao quát các công thức mô hình đích, mô hình ứng viên, hệ số Bhattacharyya và bước cập nhật vị trí ở các mục ngay sau.

\subsubsection{Histogram màu và ảnh xác suất}

Chuẩn hoá vùng chứa đối tượng về hình cầu bán kính đơn vị với tâm tại gốc toạ độ, gồm các điểm ảnh
$\{\mathbf{x}_i^{*}\}_{i=1..n}$. Gọi $b(\mathbf{x})$ là hàm ánh xạ một điểm ảnh sang chỉ số ô của
histogram màu $m$ ô. \textbf{Mô hình đích} được định nghĩa là
\begin{equation}
\hat{q}_u = C \sum_{i=1}^{n} k\!\left(\|\mathbf{x}_i^{*}\|^{2}\right)
\delta\big[b(\mathbf{x}_i^{*}) - u\big], \qquad u = 1,\dots,m,
\label{eq:target}
\end{equation}
trong đó $k$ là profile của hàm hạt nhân không gian, $\delta$ là hàm Kronecker và
$C = 1 / \sum_{i} k(\|\mathbf{x}_i^{*}\|^2)$ là hằng số chuẩn hoá bảo đảm $\sum_{u}\hat{q}_u = 1$.

Vai trò của hạt nhân không gian $k$ cần được nhấn mạnh: nó gán trọng số lớn hơn cho các điểm ảnh
gần tâm vùng đối tượng, dựa trên nhận định rằng các điểm ảnh ngoại vi kém tin cậy hơn vì dễ lẫn
với nền và thường bị che khuất trước tiên.

Tương tự, \textbf{mô hình ứng viên} tại vị trí $\mathbf{y}$ trong khung hình hiện tại là
\begin{equation}
\hat{p}_u(\mathbf{y}) = C_h \sum_{i=1}^{n_h}
k\!\left( \left\|\frac{\mathbf{y}-\mathbf{x}_i}{h}\right\|^{2}\right)
\delta\big[b(\mathbf{x}_i) - u\big],
\qquad C_h = \frac{1}{\sum_{i=1}^{n_h} k\!\left( \left\|\frac{\mathbf{y}-\mathbf{x}_i}{h}\right\|^{2}\right)},
\end{equation}
với $h$ là độ rộng dải (bandwidth) của hạt nhân.

\textbf{Phép chiếu ngược} (back-projection) là bước biến histogram thành một ảnh làm việc: mỗi
điểm ảnh của khung hình được thay bằng giá trị $\hat{q}_{b(\mathbf{x})}$, tức tần suất của màu tại
điểm đó trong mô hình đích. Kết quả là một \textit{ảnh xác suất} trong đó cường độ tại mỗi điểm
ảnh phản ánh khả năng điểm đó thuộc về đối tượng. Trong OpenCV, bước này được thực hiện bằng hàm
\texttt{cv2.calcBackProject}, và chính ảnh xác suất là đầu vào cho thủ tục tìm cực trị mật độ ở
các mục tiếp theo.

\subsubsection{Hệ số Bhattacharyya}

Mức độ tương đồng giữa mô hình đích và mô hình ứng viên được đo bằng \textbf{hệ số Bhattacharyya}
\begin{equation}
\hat{\rho}(\mathbf{y}) \;\equiv\; \rho\big[\hat{p}(\mathbf{y}), \hat{q}\big]
\;=\; \sum_{u=1}^{m} \sqrt{\hat{p}_u(\mathbf{y})\, \hat{q}_u}.
\label{eq:bhatt}
\end{equation}

Biểu thức~\eqref{eq:bhatt} có diễn giải hình học trực quan: nó là tích vô hướng của hai vector đơn
vị $(\sqrt{\hat p_1},\dots,\sqrt{\hat p_m})^\top$ và $(\sqrt{\hat q_1},\dots,\sqrt{\hat q_m})^\top$,
tức cosin của góc giữa chúng trong không gian $m$ chiều. Khoảng cách tương ứng
$d(\mathbf{y}) = \sqrt{1 - \rho[\hat{p}(\mathbf{y}), \hat{q}]}$ đã được chứng minh là một metric
hợp lệ. Nhờ đó bài toán bám được phát biểu lại thành bài toán tối ưu: tìm vị trí $\mathbf{y}$ cực
đại hoá~\eqref{eq:bhatt}.

\subsubsection{Vector mean shift}

Xét $n$ điểm dữ liệu $\mathbf{x}_i \in \mathbb{R}^d$ và hàm hạt nhân $K$ với profile $k$ thoả
$K(\mathbf{x}) = c_{k,d}\, k(\|\mathbf{x}\|^2)$. Ước lượng mật độ hạt nhân tại $\mathbf{x}$ là
\begin{equation}
\hat{f}_{h,K}(\mathbf{x}) \;=\; \frac{c_{k,d}}{n h^{d}} \sum_{i=1}^{n}
k\!\left( \left\| \frac{\mathbf{x}-\mathbf{x}_i}{h} \right\|^{2} \right).
\label{eq:kde}
\end{equation}

Lấy gradient của~\eqref{eq:kde} và đặt $g(x) = -k'(x)$, ta thu được
\begin{equation}
\nabla \hat{f}_{h,K}(\mathbf{x}) \;=\;
\frac{2 c_{k,d}}{n h^{d+2}}
\left[\sum_{i=1}^{n} g_i \right]
\underbrace{\left[ \frac{\sum_{i=1}^{n}\mathbf{x}_i\, g_i}{\sum_{i=1}^{n} g_i} - \mathbf{x} \right]}_{\textstyle \mathbf{m}_h(\mathbf{x})},
\qquad g_i = g\!\left(\left\|\tfrac{\mathbf{x}-\mathbf{x}_i}{h}\right\|^2\right).
\end{equation}

Đại lượng $\mathbf{m}_h(\mathbf{x})$ được gọi là \textbf{vector mean shift}. Do thừa số đứng trước
nó luôn dương, vector này luôn cùng hướng với gradient mật độ, tức luôn chỉ về phía mật độ tăng
nhanh nhất~\cite{comaniciu2002}. Đây là kết quả then chốt vì nó cho phép thực hiện một thủ tục leo
dốc mà không cần chọn tham số bước: độ dài bước được xác định tự động từ chính cấu trúc mật độ cục
bộ.

Hai hàm hạt nhân thường dùng là Epanechnikov với $k_E(x) = 1-x$, $0 \le x \le 1$, và Gauss với
$k_N(x) = \exp(-x/2)$. Trường hợp Epanechnikov có ý nghĩa thực tiễn đặc biệt vì khi đó $g$ là hằng
số, phép lặp suy biến thành phép lấy trung bình cộng có trọng số.

\subsubsection{Công thức cập nhật vị trí}

Khai triển Taylor bậc nhất của~\eqref{eq:bhatt} quanh các giá trị $\hat{p}_u(\mathbf{y}_0)$ tại vị
trí hiện tại $\mathbf{y}_0$ cho
\begin{equation}
\rho\big[\hat{p}(\mathbf{y}), \hat{q}\big] \approx
\frac{1}{2}\sum_{u=1}^{m} \sqrt{\hat{p}_u(\mathbf{y}_0)\hat{q}_u}
+ \frac{1}{2}\sum_{u=1}^{m} \hat{p}_u(\mathbf{y}) \sqrt{\frac{\hat{q}_u}{\hat{p}_u(\mathbf{y}_0)}} .
\end{equation}
Số hạng thứ nhất là hằng số theo $\mathbf{y}$. Thay biểu thức của $\hat p_u(\mathbf{y})$ vào số
hạng thứ hai, ta được
\begin{equation}
\rho\big[\hat{p}(\mathbf{y}), \hat{q}\big] \approx \text{const} +
\frac{C_h}{2} \sum_{i=1}^{n_h} w_i\,
k\!\left( \left\|\frac{\mathbf{y}-\mathbf{x}_i}{h}\right\|^{2}\right),
\end{equation}
với \textbf{trọng số điểm ảnh}
\begin{equation}
w_i = \sum_{u=1}^{m} \sqrt{\frac{\hat{q}_u}{\hat{p}_u(\mathbf{y}_0)}}\;
\delta\big[b(\mathbf{x}_i) - u\big].
\label{eq:weights}
\end{equation}

Ý nghĩa của~\eqref{eq:weights} rất đáng chú ý: trọng số $w_i$ lớn khi màu của điểm ảnh
$\mathbf{x}_i$ xuất hiện nhiều trong mô hình đích nhưng ít trong mô hình ứng viên hiện tại. Mỗi
điểm ảnh như vậy là một bằng chứng cho thấy cửa sổ tìm kiếm cần dịch chuyển về phía nó. Toàn bộ
thuật toán vì thế có thể hiểu như một cơ chế bỏ phiếu có trọng số giữa các điểm ảnh.

Cực đại hoá biểu thức trên chính là bài toán ước lượng mật độ có trọng số, nên lời giải là phép
lặp mean shift với công thức cập nhật vị trí
\begin{equation}
\mathbf{y}_{\text{mới}} = \frac{\displaystyle\sum_{i=1}^{n_h} \mathbf{x}_i\, w_i\,
g\!\left(\left\|\frac{\mathbf{y}_0-\mathbf{x}_i}{h}\right\|^{2}\right)}
{\displaystyle\sum_{i=1}^{n_h} w_i\,
g\!\left(\left\|\frac{\mathbf{y}_0-\mathbf{x}_i}{h}\right\|^{2}\right)}.
\label{eq:update}
\end{equation}

\begin{algorithm}[H]
\caption{Bám đối tượng bằng MeanShift trên một khung hình}
\label{alg:meanshift}
\begin{algorithmic}[1]
\Require Mô hình đích $\{\hat q_u\}$; vị trí khởi đầu $\mathbf{y}_0$ lấy từ khung hình trước;
bandwidth $h$; ngưỡng dừng $\varepsilon$; số vòng lặp tối đa $N_{\max}$
\Ensure Vị trí đối tượng ở khung hình hiện tại
\State Tính $\{\hat p_u(\mathbf{y}_0)\}$ và $\rho_0 \leftarrow \rho[\hat p(\mathbf{y}_0), \hat q]$
\State $j \leftarrow 0$
\Repeat
  \State $\mathbf{y}_{\text{trước}} \leftarrow \mathbf{y}_0$
    \Comment{lưu vị trí đầu vòng lặp để đo độ dịch chuyển}
  \State Tính trọng số $w_i$ theo~\eqref{eq:weights}
  \State Tính $\mathbf{y}_{\text{mới}}$ từ $\mathbf{y}_{\text{trước}}$ theo~\eqref{eq:update}
  \State Tính $\rho_{\text{mới}} \leftarrow \rho[\hat p(\mathbf{y}_{\text{mới}}), \hat q]$
  \While{$\rho_{\text{mới}} < \rho_0$}
     \State $\mathbf{y}_{\text{mới}} \leftarrow \tfrac{1}{2}(\mathbf{y}_{\text{trước}} + \mathbf{y}_{\text{mới}})$
     \State Cập nhật lại $\rho_{\text{mới}}$
       \Comment{lùi nửa bước, bảo đảm hệ số Bhattacharyya không giảm}
  \EndWhile
  \State $\Delta \leftarrow \|\mathbf{y}_{\text{mới}} - \mathbf{y}_{\text{trước}}\|$
    \Comment{đo độ dịch chuyển \textbf{trước khi} cập nhật $\mathbf{y}_0$}
  \State $\mathbf{y}_0 \leftarrow \mathbf{y}_{\text{mới}}$;\quad
         $\rho_0 \leftarrow \rho_{\text{mới}}$;\quad
         $j \leftarrow j + 1$
\Until{$\Delta < \varepsilon$ \textbf{hoặc} $j \ge N_{\max}$}
\State \Return $\mathbf{y}_0$
\end{algorithmic}
\end{algorithm}

Thứ tự các bước ở dòng 12 và 13 của Thuật toán~\ref{alg:meanshift} cần được tuân thủ chặt chẽ. Nếu
đo độ dịch chuyển \textit{sau} khi đã gán $\mathbf{y}_0 \leftarrow \mathbf{y}_{\text{mới}}$ thì
đại lượng thu được luôn bằng không và thuật toán sẽ dừng ngay sau vòng lặp đầu tiên. Vì vậy độ
dịch chuyển $\Delta$ phải được tính từ vị trí đầu vòng lặp $\mathbf{y}_{\text{trước}}$.

\paragraph{Tính hội tụ và độ phức tạp.} Thuật toán dùng bước lùi nửa khoảng cách khi hệ số
Bhattacharyya giảm, nhờ đó duy trì tính tăng đơn điệu của hàm tương đồng theo thủ tục trong công
trình gốc. Trong ví dụ minh hoạ của Comaniciu, Ramesh và Meer, quá trình tìm kiếm đạt nghiệm sau
bốn vòng lặp~\cite{comaniciu2003}. Độ phức tạp mỗi vòng lặp là $O(n_h)$ theo số điểm ảnh trong
cửa sổ; số vòng lặp thực tế còn phụ thuộc dữ liệu, vị trí khởi tạo và tiêu chí dừng.

\subsection{Hạn chế của MeanShift nguyên bản}

\textbf{Cửa sổ cố định.} Do bandwidth $h$ không thay đổi, thuật toán không theo kịp biến thiên tỉ
lệ của đối tượng. Giải pháp khắc phục kinh điển là chạy thêm ở $h(1\pm 10\%)$ rồi chọn giá trị cho
hệ số Bhattacharyya cao nhất, nhưng cách này thiên lệch có hệ thống về phía thu nhỏ vì cửa sổ nhỏ
hơn luôn khớp histogram tốt hơn do chứa ít điểm ảnh nền hơn. Hệ quả là cửa sổ có xu hướng co dần
cho tới khi mất mục tiêu.

\textbf{Không giữ cấu trúc không gian.} Biểu diễn bằng histogram bỏ qua hoàn toàn bố cục hình học,
do đó hai vùng ảnh có cùng phân bố màu nhưng khác hẳn nhau về hình dạng vẫn được đánh giá là tương
đương.

\textbf{Nhạy với nền cùng màu.} Khi nền chứa các thành phần màu trùng với đối tượng, trọng số
$w_i$ ở~\eqref{eq:weights} không còn phân biệt được đối tượng với nền và cửa sổ dễ bị kéo lệch.

\textbf{Không tự phát hiện lại mục tiêu.} Cơ chế tìm kiếm cục bộ về nguyên tắc không thể khôi phục
mục tiêu khi vùng đối tượng ở hai khung hình liên tiếp không giao nhau, tức khi xảy ra che khuất
toàn phần hoặc chuyển động nhanh.

\subsection{Các biến thể của MeanShift}

\begin{table}[H]
\centering
\caption{Các biến thể chính của thuật toán MeanShift}
\label{tab:variants}
\small
\begin{tabular}{@{}lp{4.1cm}p{6.4cm}@{}}
\toprule
\textbf{Thuật toán} & \textbf{Cải tiến chính} & \textbf{Cơ chế} \\
\midrule
MeanShift~\cite{comaniciu2003} & Ước lượng vị trí & Lặp mean shift trên trọng số Bhattacharyya \\
CAMShift~\cite{bradski1998} & Thích nghi kích thước và hướng & Mô-men bậc không và bậc hai của ảnh xác suất \\
CBWH~\cite{ning2012cbwh} & Giảm ảnh hưởng của màu nền & Biến đổi riêng mô hình đích theo thống kê nền \\
SOAMST~\cite{ning2012soamst} & Thích nghi tỉ lệ và góc xoay & Ellipse hiệp phương sai từ mô-men trung tâm bậc hai \\
EM-shift~\cite{zivkovic2004} & Ước lượng đồng thời vị trí và hình dạng & Lược đồ kiểu EM trên vị trí và ma trận hiệp phương sai \\
\bottomrule
\end{tabular}
\end{table}

\textbf{CAMShift} tính ảnh xác suất bằng phép chiếu ngược histogram sắc độ, sau đó điều chỉnh kích
thước cửa sổ theo mô-men bậc không $M_{00}$ đặc trưng cho diện tích phân bố, đồng thời ước lượng
hướng cùng chiều rộng và chiều cao từ các mô-men bậc hai. Đây là cơ sở của hàm
\texttt{cv2.CamShift}.

\textbf{CBWH} giảm trọng số của các thành phần màu nổi bật thuộc về nền. Đóng góp lý thuyết quan
trọng của công trình là chứng minh phương pháp BWH nguyên bản thực chất tương đương với MeanShift
chuẩn, do phép lặp mean shift bất biến với phép nhân vô hướng vào trọng số. Trên cơ sở đó nhóm tác
giả đề xuất chỉ biến đổi mô hình đích. Thực nghiệm của công trình cho thấy cách sửa này làm giảm
ảnh hưởng của nền, tăng độ chính xác định vị và hội tụ nhanh hơn so với biểu diễn đích thông thường.

\textbf{SOAMST} dùng mô-men bậc không và bậc nhất của ảnh trọng số để xác định vị trí, kết hợp hệ
số Bhattacharyya với mô-men trung tâm bậc hai để ước lượng diện tích, chiều rộng, chiều cao và góc
xoay thông qua ellipse hiệp phương sai.

\textbf{EM-shift} ước lượng đồng thời vị trí và ma trận hiệp phương sai mô tả hình dạng đối tượng
theo lược đồ kiểu EM.

% =====================================================================
\section{HỌ ĐẶC TRƯNG CỤC BỘ}
% =====================================================================

\subsection{Thuật toán SIFT}

SIFT~\cite{lowe2004} gồm bốn giai đoạn, mỗi giai đoạn xử lý một dạng biến đổi ảnh khác nhau. Cần
phát biểu chính xác về tính bất biến của thuật toán, vì cách nói gộp thường gặp dễ gây hiểu nhầm:

\begin{itemize}
  \item Đối với \textit{vị trí và tỉ lệ}, SIFT có tính \textbf{hiệp biến}: khi ảnh được tịnh tiến
  hoặc thay đổi tỉ lệ, các điểm chính dịch chuyển và thay đổi tỉ lệ tương ứng. Chính nhờ tính hiệp
  biến này mà vector mô tả tính trên vùng đã chuẩn hoá mới đạt tính bất biến.
  \item Đối với \textit{phép quay}, tính bất biến không có sẵn mà đạt được nhờ bước
  \textbf{chuẩn hoá theo hướng chính} ở giai đoạn ba.
  \item Đối với \textit{thay đổi chiếu sáng}, SIFT chỉ \textbf{bền vững} chứ không bất biến tuyệt
  đối: phép chuẩn hoá và cắt ngưỡng vector mô tả xử lý được thay đổi tương phản tuyến tính và giảm
  ảnh hưởng của hiệu ứng phi tuyến, nhưng không loại bỏ hoàn toàn.
  \item Đối với \textit{biến đổi affine mạnh hoặc thay đổi góc nhìn lớn}, SIFT \textbf{không bất
  biến hoàn toàn}; công trình gốc chỉ khẳng định độ bền vững trong một khoảng biến dạng affine và
  thay đổi góc nhìn nhất định.
\end{itemize}

\subsubsection{Không gian tỉ lệ và hàm DoG}

Không gian tỉ lệ Gauss được định nghĩa bởi $L(x,y,\sigma) = G(x,y,\sigma) * I(x,y)$. Hàm sai khác
Gauss là hiệu của hai lớp kề nhau trong không gian tỉ lệ:
\begin{equation}
D(x,y,\sigma) = \big(G(x,y,k\sigma) - G(x,y,\sigma)\big) * I(x,y) = L(x,y,k\sigma) - L(x,y,\sigma).
\end{equation}

Lựa chọn hàm DoG có cơ sở lý thuyết chặt chẽ. Từ phương trình khuếch tán
$\partial G/\partial \sigma = \sigma \nabla^2 G$, phép xấp xỉ sai phân hữu hạn cho
\begin{equation}
G(x,y,k\sigma) - G(x,y,\sigma) \approx (k-1)\,\sigma^{2} \nabla^{2} G,
\end{equation}
nghĩa là DoG xấp xỉ toán tử Laplacian chuẩn hoá theo tỉ lệ $\sigma^2 \nabla^2 G$, mà cực trị của
toán tử này cho các đặc trưng ổn định nhất qua các phép biến đổi ảnh. Ưu điểm thực tiễn là DoG
được tính bằng phép trừ đơn giản.

Không gian tỉ lệ được tổ chức thành các octave, mỗi octave chia thành $s$ khoảng với $k = 2^{1/s}$.
Tham số điển hình là $\sigma = 1{,}6$ và $s = 3$, khi đó mỗi octave sinh $s+3$ ảnh Gauss và $s+2$
ảnh DoG. Một điểm là ứng viên cực trị nếu lớn hơn hoặc nhỏ hơn toàn bộ 26 điểm lân cận, gồm 8 điểm
cùng tỉ lệ, 9 điểm ở tỉ lệ trên và 9 điểm ở tỉ lệ dưới.

Khai triển Taylor bậc hai của $D$ quanh điểm ứng viên cho vị trí và tỉ lệ ở mức dưới điểm ảnh:
\begin{equation}
\hat{\mathbf{x}} = -\left(\frac{\partial^{2} D}{\partial \mathbf{x}^{2}}\right)^{-1}
\frac{\partial D}{\partial \mathbf{x}} .
\end{equation}
Sau đó loại các điểm có độ tương phản thấp với $|D(\hat{\mathbf{x}})| < 0{,}03$, và loại các điểm
nằm trên cạnh bằng ma trận Hessian $2\times2$, chỉ giữ những điểm thoả
$\operatorname{Tr}(\mathbf{H})^{2}/\operatorname{Det}(\mathbf{H}) < (r+1)^{2}/r$ với $r = 10$. Cơ
sở của phép lọc thứ hai là điểm nằm trên cạnh có hai độ cong chính rất chênh lệch, khiến vị trí
dọc theo cạnh không xác định rõ và do đó không ổn định khi đối sánh.

\subsubsection{Gán hướng}

Quanh mỗi điểm chính, một histogram hướng gradient 36 ô được xây dựng, mỗi mẫu có trọng số bằng
tích của biên độ gradient với cửa sổ Gauss có $\sigma' = 1{,}5\sigma$. Đỉnh lớn nhất xác định
hướng chính; mọi đỉnh khác đạt từ $80\%$ đỉnh lớn nhất trở lên sinh thêm một điểm chính cùng vị
trí nhưng khác hướng. Vector mô tả sau đó được tính \textit{tương đối} so với hướng này, và đó
chính là cơ chế tạo ra tính bất biến quay.

\subsubsection{Vector mô tả 128 chiều}

Trên vùng $16\times16$ điểm ảnh quanh điểm chính, đã chuẩn hoá theo tỉ lệ và xoay theo hướng
chính, thuật toán chia thành lưới $4\times4$ ô, mỗi ô lập một histogram hướng 8 ô, thu được vector
$4\times4\times8 = 128$ chiều. Mỗi mẫu gradient được phân phối bằng nội suy tam tuyến tính nhằm
tránh hiệu ứng ranh giới khi điểm chính dịch chuyển nhẹ. Cuối cùng, vector được chuẩn hoá về độ
dài đơn vị, mọi thành phần lớn hơn $0{,}2$ bị cắt ngưỡng và vector được chuẩn hoá lại.

\subsubsection{Ratio test và RANSAC}

Việc đối sánh dùng khoảng cách Euclid giữa các vector 128 chiều, kết hợp với \textbf{kiểm định tỉ
lệ khoảng cách láng giềng gần nhất} (ratio test): một cặp khớp bị loại nếu
$d(\text{NN}_1)/d(\text{NN}_2) > 0{,}8$. Lowe~\cite{lowe2004} báo cáo ngưỡng này
\textenglish{``eliminates 90\% of the false matches while discarding less than 5\%''} số cặp khớp
đúng, trên thực nghiệm với cơ sở dữ liệu 40.000 điểm chính.

Trong bộ bám SIFT được cài đặt cho phần thực nghiệm, bước kiểm chứng hình học sử dụng
\textbf{RANSAC}: thuật toán lấy mẫu ngẫu nhiên các tập con cặp khớp và ước lượng phép biến đổi
partial affine gồm tịnh tiến, phép quay và tỉ lệ đồng nhất. Mô hình này ít bậc tự do hơn homography,
nhờ đó ổn định hơn khi mục tiêu nhỏ hoặc chỉ còn ít cặp khớp. Giả thuyết có nhiều điểm nội tại
(inlier) nhất được dùng để loại các cặp khớp sai còn sót lại sau ratio test.

\subsection{So sánh ngắn với SURF và ORB}

\begin{table}[H]
\centering
\caption{So sánh định tính giữa SIFT, SURF và ORB}
\label{tab:sift-surf-orb}
\small
\begin{tabular}{@{}lp{5.2cm}p{5.2cm}@{}}
\toprule
\textbf{Thuật toán} & \textbf{Ưu điểm} & \textbf{Hạn chế} \\
\midrule
SIFT~\cite{lowe2004} & Bền vững trước thay đổi tỉ lệ, phép quay và một khoảng biến dạng ảnh & Descriptor 128 chiều làm tăng chi phí trích xuất và đối sánh \\
SURF~\cite{bay2006} & Tăng tốc bằng bộ lọc hộp và ảnh tích phân & Là phép xấp xỉ được thiết kế ưu tiên tốc độ \\
ORB~\cite{rublee2011} & Descriptor nhị phân, đối sánh Hamming nhanh và phù hợp thời gian thực & Mục tiêu chính là hiệu quả tính toán; độ chính xác cần được đánh giá trên dữ liệu cụ thể \\
\bottomrule
\end{tabular}
\end{table}

Nhóm tác giả ORB báo cáo ORB nhanh hơn SIFT khoảng hai bậc độ lớn trong các thí nghiệm của
họ, đồng thời đạt chất lượng tương đương trong nhiều tình huống~\cite{rublee2011}. Kết quả này
không được hiểu là ORB luôn tốt hơn SIFT trên mọi dữ liệu; bảng trên chỉ nêu khác biệt thiết kế và
đánh đổi điển hình. Vì chủ đề chính của báo cáo là SIFT, SURF và ORB chỉ được dùng làm mốc tham
chiếu.

\subsection{Bám đối tượng dựa trên SIFT}

Quy trình bám dựa trên SIFT gồm bốn bước: trích xuất điểm chính trên mẫu đối tượng và trên khung
hình hiện tại; đối sánh vector mô tả kèm ratio test; ước lượng phép biến đổi tương tự, affine hoặc
homography bằng RANSAC; và suy ra thành phần tịnh tiến, tỉ lệ cùng góc xoay của đối tượng.

Về khả năng xử lý \textbf{che khuất một phần}, cần phát biểu kèm điều kiện. Phương pháp chỉ hoạt
động khi số cặp khớp đúng còn lại đủ lớn \textit{và} có phân bố hình học phù hợp. Cụ thể, việc ước
lượng phép biến đổi sẽ thất bại hoặc cho nghiệm không ổn định trong các trường hợp: số điểm nội
tại còn lại quá ít so với số bậc tự do của phép biến đổi; các điểm còn lại tập trung trong một
vùng nhỏ của đối tượng; các điểm gần như thẳng hàng, khiến hệ phương trình suy biến; hoặc tỉ lệ
điểm ngoại lai còn cao khiến RANSAC không tìm được giả thuyết đồng thuận đáng tin cậy.

Ngoài ra, phương pháp còn hai hạn chế: số điểm chính trên đối tượng nhỏ hoặc ít texture là rất ít,
thậm chí bằng không; và tập điểm chính phải được cập nhật theo thời gian, chính cơ chế cập nhật
này lại là nguồn gốc của hiện tượng trôi.

% =====================================================================
\section{SO SÁNH MEANSHIFT VÀ SIFT}
% =====================================================================

\begin{table}[H]
\centering
\caption{So sánh định tính giữa MeanShift và SIFT}
\label{tab:compare}
\small
\begin{tabularx}{\textwidth}{@{}l>{\raggedright\arraybackslash}X>{\raggedright\arraybackslash}X@{}}
\toprule
\textbf{Tiêu chí} & \textbf{MeanShift} & \textbf{SIFT} \\
\midrule
Biểu diễn đối tượng & Histogram màu toàn cục & Điểm đặc trưng cục bộ \\
Tốc độ xử lý & Cao & Thấp hơn \\
Thay đổi tỉ lệ & Yếu ở bản nguyên & Tốt \\
Phép quay & Yếu & Tốt \\
Biến dạng phi cứng & Tốt & Có thể suy giảm \\
Đối tượng ít texture & Có thể hoạt động & Dễ thất bại \\
Nền cùng màu với đối tượng & Dễ trôi & Ít phụ thuộc màu hơn \\
Che khuất một phần & Trung bình & Tốt nếu còn đủ điểm khớp \\
Phát hiện lại mục tiêu & Không có & Có thể thực hiện \\
Khả năng giải thích & Cao & Cao \\
\bottomrule
\end{tabularx}

\vspace{4pt}
\begin{minipage}{0.97\textwidth}
{\footnotesize \textit{Chú giải:} Bảng tổng hợp \textit{định tính}, được suy ra từ phân tích cơ
chế thuật toán ở Mục~2 và Mục~3, không dựa trên một thực nghiệm đối chứng chạy trên cùng một cấu
hình. Các mức đánh giá thể hiện so sánh tương đối giữa hai phương pháp trên cùng một tiêu chí.}
\end{minipage}
\end{table}

Bảng~\ref{tab:compare} cho thấy quan hệ có tính đối ngẫu giữa hai phương pháp: các tiêu chí mà
MeanShift tỏ ra yếu, gồm thay đổi tỉ lệ, phép quay và nền cùng màu, lại chính là những tiêu chí mà
SIFT xử lý tốt; ngược lại ở các tiêu chí biến dạng phi cứng, tốc độ và đối tượng ít texture.

Quan hệ đối ngẫu này không ngẫu nhiên mà bắt nguồn từ chính lựa chọn biểu diễn. Histogram toàn cục
bỏ qua cấu trúc hình học nên bất biến với biến dạng nhưng đồng thời cũng mù trước biến đổi tỉ lệ
và phép quay. Ngược lại, tập điểm đặc trưng cục bộ nắm bắt được cấu trúc hình học nên xử lý tốt
biến đổi hình học, nhưng lại đòi hỏi bề mặt có đủ chi tiết và nhạy cảm với biến dạng phi cứng làm
thay đổi quan hệ vị trí giữa các điểm.

Một khác biệt quan trọng khác nằm ở dòng cuối cùng của bảng: MeanShift không có cơ chế phát hiện
lại vì nó chỉ tìm kiếm cục bộ quanh vị trí cũ, trong khi SIFT về nguyên tắc có thể đối sánh mẫu
đối tượng với toàn bộ khung hình. Đây chính là cơ sở cho cơ chế thứ ba trong Mục~5.

% =====================================================================
\section{KẾT HỢP MEANSHIFT VÀ SIFT}
% =====================================================================

\begin{figure}[H]
\centering
\resizebox{0.98\textwidth}{!}{%
\begin{tikzpicture}[
  node distance=6mm and 8mm,
  box/.style={draw, rounded corners=2pt, align=center, font=\scriptsize, minimum height=8mm, inner sep=4pt},
  decision/.style={diamond, draw, aspect=2.2, align=center, font=\scriptsize, inner sep=1.5pt},
  arr/.style={-{Stealth[length=2.0mm]}, thick}
]
\node[box] (frame) {Khung hình\\$I_t$};
\node[box, right=8mm of frame] (ms) {MeanShift cập nhật\\hộp và độ tin cậy $\rho_t$};
\node[decision, right=9mm of ms] (trigger) {Đến chu kỳ SIFT\\hoặc $\rho_t<\theta$?};
\node[box, above right=6mm and 9mm of trigger] (keep) {Giữ kết quả\\MeanShift};
\node[box, below=9mm of trigger] (sift) {SIFT cục bộ/toàn ảnh\\ratio test + RANSAC};
\node[decision, right=9mm of sift] (valid) {Đủ cặp khớp\\và inlier?};
\node[box, below=8mm of valid] (correct) {Dùng hộp SIFT và\\đặt lại cửa sổ MeanShift};
\node[box, right=11mm of valid, yshift=9mm] (out) {Hộp bao\\$\hat{\mathbf{s}}_t$};

\draw[arr] (frame) -- (ms);
\draw[arr] (ms) -- (trigger);
\draw[arr] (trigger) -- node[above,font=\scriptsize] {Không} (keep);
\draw[arr] (trigger) -- node[left,font=\scriptsize] {Có} (sift);
\draw[arr] (sift) -- (valid);
\draw[arr] (valid) -- node[right,font=\scriptsize] {Có} (correct);
\draw[arr] (valid) |- node[pos=0.25,above,font=\scriptsize] {Không} (keep);
\draw[arr] (keep) -- (out);
\draw[arr] (correct) -| (out);
\end{tikzpicture}%
}
\caption{Luồng xử lý của bộ bám MeanShift+SIFT trong phần thực nghiệm}
\label{fig:hybrid}
\end{figure}

Trong cài đặt thực nghiệm, hai nhánh không được hợp nhất bằng trọng số trên từng pixel. MeanShift
luôn chạy trước để tạo hộp ứng viên và độ tin cậy histogram. SIFT chỉ được kích hoạt định kỳ sau
mỗi 5 khung hình hoặc khi độ tin cậy nhỏ hơn $0{,}65$. Nếu độ tin cậy còn từ $0{,}35$ trở lên,
hộp MeanShift được dùng làm gợi ý tìm kiếm; ngược lại, SIFT mở rộng tìm kiếm từ vị trí trước của
bộ bám lai. Kết quả SIFT chỉ được chấp nhận khi có ít nhất 6 cặp khớp và 4 điểm nội tại (inlier) sau RANSAC
với phép affine từng phần. Khi đó hộp SIFT đặt lại cửa sổ MeanShift; nếu không, hệ thống giữ kết quả
MeanShift. Luồng này phản ánh đúng mã thực nghiệm và tránh nhầm với phương án trọng số SIFT trên
histogram.

Ba cơ chế dưới đây giải thích vai trò của từng thành phần.

\textbf{Cơ chế 1: MeanShift bám nhanh ở mọi khung hình.} Khi mô hình màu còn đáng tin cậy,
MeanShift cung cấp kết quả cuối mà không cần chạy SIFT, nhờ đó giữ chi phí trung bình thấp.

\textbf{Cơ chế 2: SIFT hiệu chỉnh hình học.} Các cặp khớp đã qua ratio test và RANSAC cung cấp phép
biến đổi partial affine, từ đó cập nhật vị trí, tỉ lệ và góc quay của hộp bao. Hộp đã hiệu chỉnh
được chuyển lại cho MeanShift ở khung hình kế tiếp.

\textbf{Cơ chế 3: SIFT hỗ trợ tái định vị.} Khi độ tương đồng histogram giảm, SIFT mở rộng phạm vi
tìm kiếm và đối sánh với mẫu gốc. Cơ chế này chỉ phục hồi được mục tiêu khi vẫn còn đủ điểm đặc
trưng đúng; nếu không đủ điểm nội tại, bộ bám tiếp tục dùng kết quả MeanShift thay vì chấp nhận một phép
biến đổi không ổn định.

\paragraph{Vai trò giảm trôi.} Ba cơ chế trên đồng thời làm giảm hiện tượng trôi theo hai đường.
Thứ nhất, việc cập nhật bandwidth theo tỉ lệ ước lượng được ngăn cửa sổ co dần, vốn là một nguyên
nhân trôi có hệ thống. Thứ hai, các cặp khớp SIFT neo vị trí ước lượng vào mẫu đối tượng \textit{ở
khung hình đầu tiên} chứ không phải vào kết quả ở khung hình liền trước, nhờ đó cắt được vòng phản
hồi tích luỹ sai số.

\paragraph{Phương án hợp nhất theo lược đồ EM.} Zhou, Yuan và Shi kết hợp phân bố xác suất từ
nhánh SIFT và nhánh histogram màu trong một lược đồ kỳ vọng--cực đại nhằm ước lượng hợp lý cực đại
vùng khớp~\cite{zhou2009}. Công trình mô tả cơ chế hỗ trợ lẫn nhau: khi một trong hai phép đo trở
nên không ổn định, phép đo còn lại giúp duy trì quá trình bám. Báo cáo không suy diễn thêm cơ chế
ước lượng trọng số ngoài những gì nguồn này công bố rõ ràng.


% =====================================================================
\section{THỰC NGHIỆM ĐỐI CHỨNG TRÊN CÙNG MÔI TRƯỜNG}
% =====================================================================

\subsection{Mục tiêu và câu hỏi thực nghiệm}

Thực nghiệm trả lời bốn câu hỏi trong phạm vi bảy bộ bám và tám chuỗi kiểm soát.

\begin{enumerate}[label=\textbf{RQ\arabic*:},leftmargin=2.2cm]
  \item Trong cùng dữ liệu và khởi tạo, các bộ bám khác nhau như thế nào về độ chính xác,
  tỉ lệ thất bại và tốc độ?
  \item CBWH-MeanShift có giảm ảnh hưởng của vùng nền gần màu so với MeanShift nguyên bản hay không?
  \item CAMShift và SOAMST thích nghi tỉ lệ, hướng tốt hơn MeanShift cửa sổ cố định đến mức nào;
  SOAMST có ổn định hơn CAMShift hay không?
  \item MeanShift+SIFT có cải thiện độ bền đủ để bù chi phí tính toán hay không; KCF nằm ở đâu
  trong đánh đổi độ chính xác--tốc độ?
\end{enumerate}

KCF được dùng làm mốc tham chiếu ngoài hai họ chính. Trọng tâm phân tích vẫn là sự khác biệt giữa
histogram màu, đặc trưng cục bộ SIFT và phương án kết hợp; kết quả không được diễn giải như một
bảng xếp hạng chung cho toàn bộ lĩnh vực bám đối tượng.

\subsection{Các thuật toán được so sánh}

Bảy bộ bám được cài đặt bằng Python và OpenCV với cùng giao diện đầu vào--đầu ra:

\begin{table}[H]
\centering
\caption{Bảy bộ bám trong thực nghiệm đối chứng}
\label{tab:experimental-trackers}
\scriptsize
\begin{tabular}{@{}l>{\raggedright\arraybackslash}p{5.0cm}>{\raggedright\arraybackslash}p{5.8cm}@{}}
\toprule
\textbf{Bộ bám} & \textbf{Biểu diễn và cập nhật} & \textbf{Vai trò trong so sánh} \\
\midrule
MeanShift & Histogram Hue, back-projection và cửa sổ cố định & Baseline nhanh của họ histogram \\
CBWH-MeanShift & Hiệu chỉnh histogram đích bằng vùng nền; cập nhật bằng MeanShift & Đo tác dụng giảm nhiễu nền \\
CAMShift & Histogram Hue và mô-men; kích thước được làm trơn và giới hạn & Đo lợi ích của cửa sổ thích nghi \\
SOAMST~\cite{ning2012soamst} & Histogram RGB $16\times16\times16$, ảnh trọng số, mô-men và hệ số Bhattacharyya & Đo khả năng thích nghi vị trí, tỉ lệ và hướng \\
SIFT & Ratio test, RANSAC partial affine, tìm cục bộ và phát hiện lại toàn khung & Baseline đặc trưng cục bộ \\
MeanShift+SIFT & MeanShift chạy thường xuyên; SIFT hiệu chỉnh định kỳ hoặc khi độ tin cậy thấp & Đo hiệu quả của hiệu chỉnh và phát hiện lại \\
KCF~\cite{henriques2015} & Bộ lọc tương quan hạt nhân, cập nhật trực tuyến & Baseline ngoài hai họ chính \\
\bottomrule
\end{tabular}
\end{table}

Các mô hình màu được khởi tạo từ khung đầu và giữ cố định. SIFT luôn giữ mô hình gốc để phát hiện
lại, đồng thời chỉ cập nhật mô hình hoạt động sau kết quả có đủ inlier. KCF cập nhật bộ lọc trực
tuyến theo từng khung hình.

Một bộ tham số duy nhất được dùng cho cả tám chuỗi. CAMShift giới hạn thay đổi kích thước mỗi khung
ở 15\%, giữ kích thước trong khoảng $0{,}35$--$3{,}0$ lần hộp khởi tạo và làm trơn với hệ số
$0{,}45$. SOAMST dùng $\sigma=1$, $\Delta d=5$ điểm ảnh, ngưỡng hội tụ $0{,}1$, tối đa 15 vòng
lặp và giới hạn thay đổi bán trục ở 20\% mỗi khung. SIFT yêu cầu tối thiểu 6 cặp khớp, 4 inlier và
ngưỡng ratio $0{,}80$.

\subsection{Dữ liệu thực nghiệm}

Bộ mã tự sinh tám chuỗi video có ground truth ở từng khung hình. Mỗi chuỗi cô lập một dạng biến
đổi để có thể liên hệ trực tiếp kết quả với cơ chế của thuật toán.

\begin{table}[H]
\centering
\caption{Liên hệ giữa thách thức và chuỗi kiểm thử có kiểm soát}
\label{tab:synthetic-scenarios}
\scriptsize
\begin{tabular}{@{}>{\raggedright\arraybackslash}p{2.5cm}>{\ttfamily\raggedright\arraybackslash}p{4.1cm}>{\raggedright\arraybackslash}p{3.2cm}>{\raggedright\arraybackslash}p{4.3cm}@{}}
\toprule
\textbf{Thách thức} & \textnormal{\textbf{Chuỗi}} & \textbf{Biến đổi chính} & \textbf{Giả thuyết cần kiểm tra} \\
\midrule
Dịch chuyển & translation & Tịnh tiến trơn, kích thước gần cố định & Đối chiếu độ chính xác vị trí và tốc độ của các baseline \\
Thay đổi tỉ lệ & scale & Phóng to rồi thu nhỏ & So sánh cửa sổ cố định với CAMShift, SOAMST và SIFT \\
Phép quay & rotation & Góc xoay thay đổi liên tục & So sánh khả năng thích nghi hướng của CAMShift, SOAMST và SIFT \\
Che khuất & occlusion & Che một phần ở giữa chuỗi & Đánh giá khả năng duy trì và phục hồi sau che khuất \\
Nền tương đồng & background\_confusion & Vật gây nhiễu có màu gần giống mục tiêu & So sánh MeanShift với CBWH-MeanShift và KCF \\
Chuyển động nhanh & fast\_motion & Dịch chuyển lớn giữa hai khung & Kiểm tra giới hạn tìm kiếm cục bộ và tái định vị \\
Thay đổi ánh sáng & illumination & Cường độ sáng biến thiên mạnh & Đánh giá độ bền của histogram Hue và descriptor chuẩn hoá \\
Ít kết cấu & low\_texture & Mục tiêu gần đồng màu, ít chi tiết & Kiểm tra điểm yếu của SIFT và lợi thế của biểu diễn màu \\
\bottomrule
\end{tabular}
\end{table}

Mỗi chuỗi gồm 100 khung hình ở độ phân giải $480\times270$. Đây là bộ dữ liệu \textit{kiểm soát},
phù hợp để giải thích nguyên nhân thành công hoặc thất bại, nhưng không đại diện cho video tự nhiên
quy mô lớn.

Hình~\ref{fig:data-a} và Hình~\ref{fig:data-b} minh hoạ tám chuỗi cùng hộp bao ground truth, giúp
đối chiếu trực tiếp đặc điểm dữ liệu với kết quả định lượng.

\begin{figure}[H]
\centering
\includegraphics[width=0.96\textwidth]{figures/scenario_examples_a.png}
\caption{Nhóm biến đổi hình học và che khuất: \scenario{translation}, \scenario{scale},
\scenario{rotation} và \scenario{occlusion}.}
\label{fig:data-a}
\end{figure}

\begin{figure}[H]
\centering
\includegraphics[width=0.96\textwidth]{figures/scenario_examples_b.png}
\caption{Các chuỗi còn lại: \scenario{background\_confusion}, \scenario{fast\_motion},
\scenario{illumination} và \scenario{low\_texture}.}
\label{fig:data-b}
\end{figure}

\subsection{Độ đo đánh giá}

Với hộp bao dự đoán $r_t$ và ground truth $r_t^{*}$, tỉ lệ chồng lấp được tính bởi
\begin{equation}
\operatorname{IoU}_t = \frac{|r_t \cap r_t^{*}|}{|r_t \cup r_t^{*}|}.
\end{equation}

\textbf{Mean IoU} là IoU trung bình trên toàn bộ khung hình. \textbf{Success plot} biểu diễn tỉ lệ
khung hình có $\operatorname{IoU}_t\ge\tau$ với $\tau\in[0,1]$; diện tích dưới đường cong là
\textbf{AUC}.

Sai số vị trí tâm được tính bởi
\begin{equation}
\operatorname{CLE}_t = \left\|\mathbf{c}_t-\mathbf{c}_t^{*}\right\|_2,
\end{equation}
trong đó $\mathbf{c}_t$ và $\mathbf{c}_t^{*}$ là tâm hộp dự đoán và ground truth.
\textbf{Precision@20} là tỉ lệ khung hình có $\operatorname{CLE}_t\le20$ điểm ảnh; một khung hình
được tính là \textbf{thất bại} khi $\operatorname{IoU}_t<0{,}1$.

Tám chuỗi có cùng độ dài nên các chỉ số được tổng hợp trực tiếp trên toàn bộ khung hình. Tốc độ chỉ
đo bước \texttt{update}. Mỗi cặp thuật toán--chuỗi chạy ba lần để lấy trung vị FPS; dự đoán của lượt
đầu được dùng cho các chỉ số độ chính xác và video minh hoạ.

\subsection{Quy trình bảo đảm tính công bằng và tái lập}

Bảng~\ref{tab:fairness} tóm tắt giao thức cố định áp dụng cho mọi bộ bám.

\begin{table}[H]
\centering
\caption{Quy trình bảo đảm tính công bằng và khả năng tái lập}
\label{tab:fairness}
\small
\begin{tabular}{@{}p{3.7cm}p{6.0cm}p{4.0cm}@{}}
\toprule
\textbf{Nguyên tắc} & \textbf{Cách thực hiện} & \textbf{Mục đích} \\
\midrule
Khởi tạo đồng nhất & Mọi tracker nhận cùng hộp bao ground truth ở khung đầu. & Loại bỏ sai lệch do khởi tạo. \\
Tham số cố định & Một bộ tham số dùng cho cả tám chuỗi; không tinh chỉnh riêng sau khi xem kết quả. & So sánh cơ chế thay vì công sức tuning. \\
Dữ liệu và RNG cố định & Các chuỗi dùng seed 100--107; OpenCV dùng seed 20260801. & Cho phép tái tạo cùng dữ liệu và các bước ngẫu nhiên của RANSAC. \\
Điều kiện đo tốc độ thống nhất & Chỉ dùng CPU và giới hạn OpenCV ở một luồng. & Tránh lợi thế do số luồng khác nhau. \\
Tách độ chính xác và tốc độ & Lượt đầu dùng cho độ chính xác; ba lần lặp chỉ dùng để lấy trung vị FPS. & Không trộn biến thiên thời gian với kết quả định vị. \\
Lệnh tái lập & \texttt{python experiment/run\_all.py --repeats 3 --save-videos} & Sinh lại CSV, biểu đồ, bảng và video minh hoạ. \\
\bottomrule
\end{tabular}
\end{table}

\IfFileExists{results/hardware.tex}{
  Bảng~\ref{tab:hardware} ghi lại môi trường của lần chạy chính thức. Toàn bộ phép đo sử dụng CPU;
  GPU không tham gia quá trình thực nghiệm.
  \input{results/hardware.tex}
}{
  \begin{center}
  \fbox{\parbox{0.92\textwidth}{\textbf{Chưa có cấu hình thực nghiệm.}
  Chạy \texttt{python experiment/run\_all.py --repeats 3 --save-videos} để sinh
  \texttt{results/hardware.tex}.}}
  \end{center}
}

\subsection{Kết quả và phân tích}

\IfFileExists{results/soamst_complete.flag}{
  \input{results/generated_results.tex}
}{
  \begin{center}
  \fbox{\parbox{0.92\textwidth}{\textbf{Kết quả hiện có chưa bao gồm SOAMST.}
  Chạy \texttt{python experiment/run\_all.py --only-soamst --repeats 3 --save-videos} để bổ sung
  SOAMST, hoặc chạy lại toàn bộ bằng \texttt{python experiment/run\_all.py --repeats 3 --save-videos}.
  Sau đó biên dịch báo cáo bằng XeLaTeX hai lần.}}
  \end{center}
}

\subsection{Các dạng thất bại điển hình}

Bảng~\ref{tab:failure-modes} tổng hợp các lỗi có thể quan sát và liên hệ chúng với cơ chế của từng
bộ bám. Mục tiêu của bảng không phải gán một lỗi duy nhất cho mỗi thuật toán, mà chỉ ra tín hiệu
nào cần kiểm tra khi kết quả suy giảm.

\begin{table}[H]
\centering
\caption{Các dạng thất bại điển hình và nguyên nhân}
\label{tab:failure-modes}
\footnotesize
\begin{tabular}{@{}>{\raggedright\arraybackslash}p{3.3cm}>{\raggedright\arraybackslash}p{4.2cm}>{\raggedright\arraybackslash}p{6.1cm}@{}}
\toprule
\textbf{Dạng thất bại} & \textbf{Bộ bám dễ gặp} & \textbf{Nguyên nhân chính} \\
\midrule
Trôi sang vùng nền cùng màu & MeanShift, CAMShift, SOAMST & Histogram hoặc mô-men phản hồi mạnh với vật gây nhiễu có phân bố màu gần mục tiêu \\
Hộp bao phình hoặc co sai & CAMShift, SOAMST & Mô-men bậc hai bị ảnh hưởng bởi nền, che khuất hoặc vùng xác suất phân tán \\
Không đủ điểm đặc trưng & SIFT, MeanShift+SIFT & Mục tiêu nhỏ, ít kết cấu hoặc bị làm mờ nên không tạo đủ vector mô tả ổn định \\
Đối sánh hình học sai & SIFT, MeanShift+SIFT & Nhiều điểm ngoại lai hoặc số điểm nội tại không đủ để ước lượng phép affine từng phần đáng tin cậy \\
Không theo kịp chuyển động & MeanShift, KCF & Mục tiêu rời khỏi vùng tìm kiếm cục bộ giữa hai khung liên tiếp \\
Không phục hồi sau che khuất dài & Các bộ bám cục bộ; bộ lai khi thiếu điểm nội tại & Không có bộ phát hiện toàn cục độc lập; SIFT chỉ tái định vị được khi mẫu còn đủ điểm khớp \\
\bottomrule
\end{tabular}
\end{table}

Phân tích lỗi cho thấy độ chính xác không chỉ phụ thuộc vào thuật toán cập nhật vị trí mà còn phụ
thuộc trực tiếp vào biểu diễn mục tiêu. Histogram màu cần độ tương phản với nền; mô-men cần vùng
xác suất tập trung; còn SIFT cần đủ cấu trúc gradient và điểm nội tại. Vì vậy, một bộ bám tốt trên một
chuỗi chưa chắc giữ lợi thế khi điều kiện quan sát thay đổi.

\subsection{Giới hạn của thực nghiệm}

Kết quả phụ thuộc vào các cài đặt cụ thể trong bộ mã. MeanShift, CBWH và CAMShift dùng histogram
Hue một chiều; SOAMST dùng histogram RGB $16\times16\times16$. SOAMST nguyên bản biểu diễn mục
tiêu bằng ellipse có hướng, trong khi giao thức này chỉ cung cấp hộp chữ nhật song song trục ở khung
đầu; vì vậy chuỗi \scenario{rotation} có thể bất lợi cho SOAMST. SIFT dùng mô hình gốc và mô hình
hoạt động cập nhật thận trọng; KCF dùng đặc trưng xám với Gaussian kernel. Thay đổi không gian màu,
số ô histogram hoặc các tham số của SOAMST, SIFT và KCF có thể làm thay đổi kết quả.

Dữ liệu tổng hợp giúp cô lập từng biến đổi nhưng không bao phủ đầy đủ nhiễu, biến dạng và độ đa
dạng của video tự nhiên. Do không đánh giá trên benchmark video thực, các số liệu chỉ có ý nghĩa
trong phạm vi tám chuỗi kiểm soát đã mô tả.

% =====================================================================
\section{KẾT LUẬN}
% =====================================================================

Báo cáo tập trung vào ba nhóm có quan hệ trực tiếp: họ MeanShift, họ đặc trưng cục bộ với SIFT
làm trọng tâm, và phương án kết hợp MeanShift+SIFT; KCF được bổ sung như baseline bộ lọc tương
quan. Cách thu hẹp phạm vi này cho phép so sánh các thuật toán theo cùng một câu hỏi: biểu diễn đối
tượng bằng phân bố màu toàn cục hay bằng tập điểm đặc trưng cục bộ sẽ dẫn tới ưu, nhược điểm gì
trong bám video.

Về mặt cơ chế, MeanShift có cấu trúc đơn giản và chi phí thấp; CBWH bổ sung thông tin nền vào mô
hình đích; CAMShift thích nghi kích thước và hướng từ mô-men ảnh xác suất. SOAMST sử dụng ảnh trọng
số phụ thuộc đồng thời mô hình đích và mô hình ứng viên, rồi kết hợp hệ số Bhattacharyya với mô-men
để ước lượng diện tích, hai trục và hướng. SIFT cung cấp khả năng đối sánh hình học và phát hiện lại
khi còn đủ điểm đặc trưng. MeanShift+SIFT tận dụng hai nguồn thông tin bổ trợ, trong khi KCF cung cấp
mốc tham chiếu của hướng học bộ lọc tương quan trực tuyến.

Thực nghiệm trên tám chuỗi tổng hợp cô lập từng yếu tố để giải thích cơ chế thành công hoặc thất
bại: \scenario{scale} và \scenario{rotation} đối chiếu trực tiếp CAMShift với SOAMST và các bộ bám khác; \scenario{low\_texture}
kiểm tra điểm yếu của SIFT; \scenario{background\_confusion} so sánh MeanShift với CBWH;
\scenario{fast\_motion} kiểm tra giới hạn của tìm kiếm cục bộ; và \scenario{occlusion} kiểm tra khả
năng phục hồi của nhánh SIFT. Các kết luận định lượng cuối cùng dựa trên bảng và biểu đồ được sinh
từ chính lần chạy trên máy của nhóm. Do không sử dụng benchmark video tự nhiên, kết quả chỉ được
hiểu trong phạm vi bộ dữ liệu kiểm soát và không đại diện cho thứ hạng tổng quát ngoài thực nghiệm.

Trong lần chạy chính thức, MeanShift+SIFT đạt kết quả tổng thể cao nhất với AUC 0,829,
Precision@20 93,8\% và tỉ lệ thất bại 1,4\%. SOAMST đạt AUC 0,501 và 450,3 FPS. Trên
\scenario{scale}, SOAMST cải thiện Mean IoU từ 0,448 của MeanShift lên 0,556, nhưng vẫn thấp hơn
CAMShift (0,715) và SIFT (0,867). Trên \scenario{rotation}, SOAMST đạt 0,618, thấp hơn CAMShift
(0,755). Vì vậy, kết quả trả lời RQ3 theo hướng: cơ chế thích nghi của SOAMST có ích so với cửa sổ
cố định trong bài toán thay đổi tỉ lệ, nhưng cài đặt thực nghiệm này chưa cho thấy SOAMST ổn định
hơn CAMShift về tổng thể hoặc khi góc quay thay đổi. Điểm nổi bật riêng của SOAMST là Mean IoU
0,985 trên \scenario{low\_texture}, cao nhất trong bảy bộ bám, phù hợp với lợi thế của biểu diễn
phân bố màu khi đối tượng thiếu điểm đặc trưng cục bộ.

Một ưu điểm chung của các phương pháp khảo sát là khả năng giải thích. Sai lệch có thể truy về
histogram và độ tương đồng Bhattacharyya, mô-men hình học, số lượng cặp khớp và điểm nội tại hoặc đỉnh đáp ứng
của bộ lọc tương quan. Các bộ bám này không cần dữ liệu huấn luyện và có thể chạy trên CPU, nên
vẫn phù hợp cho mục đích học tập, kiểm chứng thuật toán và các hệ thống hạn chế tài nguyên.

Giá trị chính của báo cáo không nằm ở việc đề xuất thuật toán mới, mà ở ba điểm: trình bày cơ sở
toán học của MeanShift và SIFT; chỉ ra quan hệ bổ trợ giữa hai cách biểu diễn; và cung cấp một quy
trình thực nghiệm nhỏ gọn, công bằng và có khả năng tái lập để người đọc tự kiểm chứng các nhận
định đó.

% =====================================================================
\begin{thebibliography}{99}
\addcontentsline{toc}{section}{TÀI LIỆU THAM KHẢO}
\footnotesize
\setlength{\itemsep}{0.15em}

\bibitem{smeulders2014}
A. W. M. Smeulders, D. M. Chu, R. Cucchiara, S. Calderara, A. Dehghan, and M. Shah,
``Visual tracking: An experimental survey,'' \emph{IEEE Transactions on Pattern Analysis and
Machine Intelligence}, vol. 36, no. 7, pp. 1442--1468, 2014.
DOI: 10.1109/TPAMI.2013.230.

\bibitem{wu2015}
Y. Wu, J. Lim, and M.-H. Yang, ``Object tracking benchmark,'' \emph{IEEE Transactions on Pattern
Analysis and Machine Intelligence}, vol. 37, no. 9, pp. 1834--1848, 2015.
DOI: 10.1109/TPAMI.2014.2388226.

\bibitem{comaniciu2002}
D. Comaniciu and P. Meer, ``Mean shift: A robust approach toward feature space analysis,''
\emph{IEEE Transactions on Pattern Analysis and Machine Intelligence}, vol. 24, no. 5,
pp. 603--619, 2002. DOI: 10.1109/34.1000236.

\bibitem{comaniciu2003}
D. Comaniciu, V. Ramesh, and P. Meer, ``Kernel-based object tracking,''
\emph{IEEE Transactions on Pattern Analysis and Machine Intelligence}, vol. 25, no. 5,
pp. 564--577, 2003. DOI: 10.1109/TPAMI.2003.1195991.

\bibitem{bradski1998}
G. R. Bradski, ``Computer vision face tracking for use in a perceptual user interface,''
\emph{Intel Technology Journal}, Q2, pp. 1--15, 1998.

\bibitem{ning2012cbwh}
J. Ning, L. Zhang, D. Zhang, and C. Wu, ``Robust mean-shift tracking with corrected
background-weighted histogram,'' \emph{IET Computer Vision}, vol. 6, no. 1, pp. 62--69, 2012.
DOI: 10.1049/iet-cvi.2009.0075.

\bibitem{ning2012soamst}
J. Ning, L. Zhang, D. Zhang, and C. Wu, ``Scale and orientation adaptive mean shift tracking,''
\emph{IET Computer Vision}, vol. 6, no. 1, pp. 52--61, 2012.
DOI: 10.1049/iet-cvi.2010.0112.

\bibitem{zivkovic2004}
Z. Zivkovic and B. Kr\"ose, ``An EM-like algorithm for color-histogram-based object tracking,''
in \emph{Proceedings of the IEEE Conference on Computer Vision and Pattern Recognition},
vol. 1, pp. 798--803, 2004. DOI: 10.1109/CVPR.2004.1315113.

\bibitem{lowe2004}
D. G. Lowe, ``Distinctive image features from scale-invariant keypoints,''
\emph{International Journal of Computer Vision}, vol. 60, no. 2, pp. 91--110, 2004.
DOI: 10.1023/B:VISI.0000029664.99615.94.

\bibitem{bay2006}
H. Bay, T. Tuytelaars, and L. Van Gool, ``SURF: Speeded up robust features,'' in
\emph{Computer Vision--ECCV 2006}, Lecture Notes in Computer Science, vol. 3951,
pp. 404--417, 2006. DOI: 10.1007/11744023\_32.

\bibitem{rublee2011}
E. Rublee, V. Rabaud, K. Konolige, and G. Bradski, ``ORB: An efficient alternative to SIFT or
SURF,'' in \emph{Proceedings of the IEEE International Conference on Computer Vision},
pp. 2564--2571, 2011. DOI: 10.1109/ICCV.2011.6126544.

\bibitem{henriques2015}
J. F. Henriques, R. Caseiro, P. Martins, and J. Batista, ``High-speed tracking with kernelized
correlation filters,'' \emph{IEEE Transactions on Pattern Analysis and Machine Intelligence},
vol. 37, no. 3, pp. 583--596, 2015. DOI: 10.1109/TPAMI.2014.2345390.

\bibitem{zhou2009}
H. Zhou, Y. Yuan, and C. Shi, ``Object tracking using SIFT features and mean shift,''
\emph{Computer Vision and Image Understanding}, vol. 113, no. 3, pp. 345--352, 2009.
DOI: 10.1016/j.cviu.2008.08.006.

\end{thebibliography}

\end{document}
